% capitulo6_implementacion.tex
% Capitulo 6: Implementacion y entorno experimental
%
% Autor: Leon Elliott Fuller
% Fecha: Mayo 2026

\chapter{Implementación y entorno experimental}
\label{ch:implementacion}

\section{Visión general}
\label{sec:impl-vision}

El objetivo de la parte práctica de este trabajo es construir una suite de benchmarking que permita comparar, de forma reproducible y cuantificable, el rendimiento de RSA y ECC a lo largo de tres ejes ortogonales:

\begin{enumerate}
    \item \textbf{RSA frente a ECC}, a niveles de seguridad equivalentes según NIST SP~800-57~\cite{nist-sp800-57}.
    \item \textbf{Coordenadas afines frente a Jacobianas}, midiendo el impacto real de eliminar inversiones modulares.
    \item \textbf{Cuerpo primo $\mathbb{F}_p$ frente a cuerpo binario $\mathbb{F}_{2^m}$}, comparando la aritmética subyacente.
\end{enumerate}

La implementación se ha desarrollado íntegramente en C++17, utilizando la biblioteca NTL~\cite{shoup-ntl} para aritmética de precisión arbitraria y GMP~\cite{granlund-gmp} como motor de cálculo subyacente. Todo el entorno se ejecuta dentro de un contenedor Docker para garantizar la reproducibilidad de los resultados.

En las secciones siguientes se describe la arquitectura del software, las decisiones de diseño tomadas, el entorno de ejecución, y los procedimientos de validación aplicados al generador de números aleatorios.


% ============================================================================
\section{Arquitectura del software}
\label{sec:impl-arquitectura}

\subsection{Estructura modular}
\label{subsec:impl-modulos}

El proyecto sigue una arquitectura modular en la que cada componente criptográfico reside en un par de ficheros (cabecera \texttt{.hpp} e implementación \texttt{.cpp}). La figura~\ref{fig:arquitectura} muestra el diagrama de dependencias.

\begin{figure}[htbp]
\centering
\begin{tikzpicture}[
    every node/.style={font=\small\ttfamily, draw, rounded corners, minimum height=0.8cm, minimum width=2cm, align=center, fill=blue!5},
    arrow/.style={->, >=latex, thick},
    node distance=0.4cm and 0.6cm
]
    % Capa 0: punto de entrada
    \node (main) [fill=orange!15] {main.cpp};

    % Capa 1: modulos criptograficos
    \node (rsa) [below left=1.2cm and 1.5cm of main] {rsa.\{hpp,cpp\}};
    \node (ecc) [below=1.2cm of main] {ecc.\{hpp,cpp\}};
    \node (eccb) [below right=1.2cm and -0.4cm of main] {ecc\_binary.\{hpp,cpp\}};
    \node (sha) [right=0.6cm of eccb, fill=blue!5] {sha256.\{hpp,cpp\}};

    % Capa 2: utilidades base
    \node (rng) [below=1.4cm of ecc, fill=green!10] {rng.\{hpp,cpp\}};
    \node (common) [right=1.0cm of rng, fill=green!10] {common.hpp};

    % Aristas desde main
    \draw [arrow] (main) -- (rsa);
    \draw [arrow] (main) -- (ecc);
    \draw [arrow] (main) -- (eccb);
    \draw [arrow] (main.east) to[bend left=15] (sha);

    % Aristas a la capa de utilidades
    \draw [arrow] (rsa) -- (rng);
    \draw [arrow] (rsa.south) -- (common.north west);
    \draw [arrow] (ecc) -- (rng);
    \draw [arrow] (ecc.south east) -- (common.north);
    \draw [arrow] (eccb.south west) -- (common.north east);
    \draw [arrow] (eccb) -- (rng.north east);
    \draw [arrow] (sha.south) -- (common.north east);
\end{tikzpicture}
\caption{Diagrama de dependencias entre los módulos del proyecto. \texttt{main.cpp} actúa como punto de entrada y orquesta los módulos criptográficos (RSA, ECC sobre $\mathbb{F}_p$ y $\mathbb{F}_{2^m}$, SHA-256), todos los cuales se apoyan en el generador de números aleatorios \texttt{rng} y los tipos comunes definidos en \texttt{common.hpp}.}
\label{fig:arquitectura}
\end{figure}

Los módulos del proyecto son los siguientes:

\begin{description}
    \item[\texttt{common.hpp}] Define los tipos base compartidos por todos los módulos. El alias \texttt{BigInt = NTL::ZZ} proporciona enteros de precisión arbitraria. Incluye constantes globales (tamaños de clave RSA por defecto, número de iteraciones Miller-Rabin, curva ECC por defecto) y una excepción personalizada \texttt{CryptoException}.
    
    \item[\texttt{rng.hpp / rng.cpp}] Implementa el generador de números aleatorios. Define una interfaz abstracta \texttt{RNG} con métodos \texttt{random\_bnd}, \texttt{random\_range}, \texttt{random\_len}, \texttt{random\_bits} y \texttt{random\_prime}. La implementación concreta \texttt{NTLRNG} delega en las funciones criptográficas de NTL y soporta semillas fijas para reproducibilidad.
    
    \item[\texttt{rsa.hpp / rsa.cpp}] Implementa RSA completo: generación de claves (con parámetros CRT precalculados), cifrado, descifrado (con y sin optimización CRT), firma y verificación. La estructura \texttt{RSAPrivateKey} almacena $p$, $q$, $d_p = d \bmod (p-1)$, $d_q = d \bmod (q-1)$ y $q^{-1} \bmod p$ para descifrado CRT.
    
    \item[\texttt{ecc.hpp / ecc.cpp}] Implementa la criptografía de curvas elípticas sobre cuerpos primos. Incluye dos representaciones de puntos: coordenadas afines (\texttt{ECPoint}) y coordenadas Jacobianas (\texttt{JacobianPoint}). Soporta tres curvas estándar: NIST P-256, NIST P-384 y secp256k1. Implementa generación de claves, ECDH, y ECDSA (firma y verificación). Todas las operaciones aceptan un parámetro \texttt{use\_jacobian} que permite seleccionar la representación en tiempo de ejecución.
    
    \item[\texttt{ecc\_binary.hpp / ecc\_binary.cpp}] Implementa curvas elípticas sobre cuerpos binarios $\mathbb{F}_{2^m}$ utilizando \texttt{NTL::GF2E} para la aritmética de polinomios. Soporta cinco curvas estándar SEC~2: sect163k1, sect233k1, sect283k1, sect233r1 y sect283r1. Incluye generación de claves y ECDH sobre cuerpos binarios.
    
    \item[\texttt{sha256.hpp / sha256.cpp}] Implementación completa de SHA-256 siguiendo FIPS~PUB~180-4~\cite{nist-fips180-4}. Se implementó desde cero (sin delegar en OpenSSL) por dos razones: integridad académica (comprender y documentar cada paso del algoritmo) y control total sobre el benchmarking (evitar que variaciones en la versión de OpenSSL afecten las mediciones).
    
    \item[\texttt{main.cpp}] Motor de benchmarking con cinco modos de ejecución, parseo de argumentos CLI mediante \texttt{getopt}, y generación de resultados en formato CSV.
\end{description}


\subsection{Interfaz de línea de comandos}
\label{subsec:impl-cli}

El binario compilado acepta los siguientes parámetros:

\begin{verbatim}
./bin/bench [opciones]
  -a ALGO    RSA | ECC | ECCJ | BIN | CMP   (default: RSA)
  -b BITS    Tamano de clave RSA              (default: 2048)
  -c CURVE   Nombre de curva ECC              (default: secp256k1)
  -i ITERS   Numero de iteraciones            (default: 50)
  -s SEED    fixed | random                   (default: fixed)
  -r FILE    Fichero CSV de datos brutos
  -v         Modo verboso (progreso a stderr)
\end{verbatim}

Los cinco modos de algoritmo corresponden a los tres ejes de comparación:

\begin{table}[htbp]
\centering
\begin{tabular}{lll}
\hline
\textbf{Modo} & \textbf{Descripción} & \textbf{Etiqueta CSV} \\
\hline
\texttt{RSA}  & RSA con tamaños 1024--4096 bits        & \texttt{RSA} \\
\texttt{ECC}  & ECC sobre $\mathbb{F}_p$ (coordenadas afines)   & \texttt{ECC} \\
\texttt{ECCJ} & ECC sobre $\mathbb{F}_p$ (coordenadas Jacobianas) & \texttt{ECC\_JACOBIAN} \\
\texttt{BIN}  & ECC sobre $\mathbb{F}_{2^m}$            & \texttt{ECC\_BINARY} \\
\texttt{CMP}  & Comparación completa (los cuatro anteriores) & Todas \\
\hline
\end{tabular}
\caption{Modos de ejecución del motor de benchmarking.}
\label{tab:cli-modes}
\end{table}

El modo \texttt{CMP} ejecuta automáticamente los cuatro modos restantes y genera aproximadamente 72 benchmarks individuales: 4 tamaños RSA $\times$ 6 operaciones, 3 curvas primas $\times$ 6 operaciones (afines), 3 curvas primas $\times$ 5 operaciones (Jacobianas), y 5 curvas binarias $\times$ 3 operaciones.


% ============================================================================
\section{Decisiones de diseño}
\label{sec:impl-decisiones}

\subsection{Progresión pedagógica: afines, después Jacobianas}
\label{subsec:impl-progresion}

Una decisión fundamental del proyecto fue implementar primero la aritmética de curvas en coordenadas afines y, solo después de validar su corrección, anadir las coordenadas Jacobianas como optimización explícita. Esta decisión responde a un principio de ingenieria de software y a un criterio académico:

\begin{description}
    \item[Claridad matemática.] La ecuación de Weierstrass $y^2 = x^3 + ax + b$ se verifica directamente sustituyendo $(x, y)$. En coordenadas Jacobianas, el mismo punto tiene infinitas representaciones equivalentes $(X, Y, Z) \sim (\lambda^2 X, \lambda^3 Y, \lambda Z)$ para cualquier $\lambda \neq 0$, lo que dificulta la comprobación visual.
    
    \item[Verificabilidad.] Después de cada operación (suma, duplicación), el resultado se puede comprobar con \texttt{is\_on\_curve()}. En coordenadas Jacobianas hay que convertir a afines antes de comparar con vectores de test conocidos.
    
    \item[Corrección primero.] El principio seguido fue ``primero que funcione correctamente, luego optimizar''. Las coordenadas afines proporcionan la base de referencia contra la que validar cualquier optimización.
    
    \item[Valor académico.] La progresión afines $\to$ Jacobianas es en sí misma un resultado medible del TFG: el speedup observado (1.7--1.9$\times$ en la práctica, frente al teórico 8.3$\times$) resulta más modesto que la estimación teórica, por las razones que se analizan en el capítulo de resultados.
\end{description}


\subsection{SHA-256 desde cero}
\label{subsec:impl-sha256}

La implementación de SHA-256 se realizó siguiendo paso a paso el estándar FIPS~PUB~180-4, sin utilizar la implementación de OpenSSL. Las razones son:

\begin{itemize}
    \item \textbf{Integridad académica:} En un TFG sobre criptografía, es importante demostrar la comprension de los algoritmos, no solo invocar bibliotecas de terceros.
    \item \textbf{Control del benchmarking:} Al implementar SHA-256 nosotros mismos, el hash no introduce una variable externa que pueda variar entre versiones de OpenSSL o configuraciones de compilación.
    \item \textbf{Coherencia:} El mismo principio se aplicó a RSA, ECC y el RNG: toda la criptografía del proyecto esta implementada sobre una única base (NTL/GMP), sin dependencias externas salvo las aritméticas.
\end{itemize}


\subsection{Uso de curvas estandarizadas}
\label{subsec:impl-curvas-estandar}

En lugar de generar curvas elípticas propias, el proyecto utiliza exclusivamente curvas estandarizadas publicadas por el NIST~\cite{nist-fips186-4} y SECG~\cite{sec2-v2}. Los parámetros del dominio $(p, a, b, G, n, h)$ están codificados directamente en \texttt{get\_curve\_params()} y \texttt{get\_binary\_curve\_params()}.

Esta decisión se justifica por:
\begin{itemize}
    \item La generación de curvas seguras requiere algoritmos de conteo de puntos (Schoof--Elkies--Atkin) que están fuera del alcance de este trabajo.
    \item Las curvas NIST y SEC~2 han sido auditadas extensamente por la comunidad criptográfica.
    \item Usar curvas estándar permite comparar nuestros resultados con benchmarks de referencia (por ejemplo, OpenSSL).
\end{itemize}

No obstante, se implementa una validación parcial de los parámetros (sección~\ref{subsec:impl-validacion-params}) para verificar las propiedades que son computacionalmente viables.


\subsection{Validación de parámetros del dominio}
\label{subsec:impl-validacion-params}

La función \texttt{CurveParams::validate()} implementa las siguientes comprobaciones para curvas sobre $\mathbb{F}_p$:

\begin{enumerate}
    \item $p > 3$.
    \item Discriminante no nulo: $4a^3 + 27b^2 \not\equiv 0 \pmod{p}$.
    \item El punto generador $G$ pertenece a la curva: $G_y^2 \equiv G_x^3 + aG_x + b \pmod{p}$.
    \item $n > 0$ y $h > 0$.
\end{enumerate}

De forma analoga, \texttt{BinaryCurveParams::validate()} comprueba que $\deg(f) = m$, que $b \neq 0$ (discriminante en característica~2), y que $n, h > 0$.

Las verificaciones que \emph{no} se implementan ---primalidad de $n$, $nG = \mathcal{O}$, curva no anómala, grado de inmersión alto--- se describen como requisitos teóricos en la sección~\ref{subsec:ecc-parametros-dominio}. Su omisión esta justificada por el uso exclusivo de curvas pre-auditadas (véase la discusión en la sección~\ref{subsubsec:verificacion-params}).


% ============================================================================
\section{Herramientas y entorno de ejecución}
\label{sec:impl-entorno}

\subsection{Lenguaje y bibliotecas}
\label{subsec:impl-lenguaje}

El cuadro~\ref{tab:herramientas} resume las herramientas principales del proyecto.

\begin{table}[htbp]
\centering
\begin{tabular}{lll}
\hline
\textbf{Herramienta} & \textbf{Versión} & \textbf{Función} \\
\hline
C++ & C++17 (ISO/IEC 14882:2017) & Lenguaje de implementación \\
NTL & $\geq$ 11.5 & Aritmética de precisión arbitraria, $\mathbb{F}_p$, $\mathbb{F}_{2^m}$ \\
GMP & $\geq$ 6.2 & Motor aritmético subyacente de NTL \\
g++ & $\geq$ 9.0 & Compilador (flags: \texttt{-std=c++17 -O2}) \\
Docker & $\geq$ 20.0 & Contenedorizacion del entorno \\
Python 3 & $\geq$ 3.8 & Visualización y análisis estadístico \\
matplotlib & $\geq$ 3.5 & Generación de gráficas \\
Pillow & $\geq$ 9.0 & Composición de collages \\
\hline
\end{tabular}
\caption{Herramientas y bibliotecas del proyecto.}
\label{tab:herramientas}
\end{table}

\paragraph{Eleccion de NTL.}
Se eligió NTL (Number Theory Library) de Victor Shoup~\cite{shoup-ntl} como base aritmética por varias razones: proporciona tipos nativos para enteros de precisión arbitraria (\texttt{ZZ}), aritmética modular (\texttt{ZZ\_p}), y cuerpos binarios de extensión (\texttt{GF2E}); ofrece funciones criptográficas de bajo nivel (exponenciación modular, test de primalidad Miller-Rabin, generación de primos); y se construye sobre GMP, heredando sus optimizaciones a nivel de ensamblador para las operaciones aritméticas básicas.

\paragraph{Nivel de optimización.}
Se compila con \texttt{-O2} en lugar de \texttt{-O3} para mantener un equilibrio entre rendimiento y predictibilidad de los tiempos de ejecución. El nivel \texttt{-O3} puede activar vectorización y desenrollado de bucles que introduzcan variabilidad en las mediciones.


\subsection{Contenedorizacion con Docker}
\label{subsec:impl-docker}

Para garantizar la reproducibilidad de los resultados, el proyecto incluye un \texttt{Dockerfile} que define un entorno de compilación y ejecución completo basado en Ubuntu~24.04. El contenedor instala NTL, GMP, Python~3 con las bibliotecas de visualización, compila el proyecto, y permite ejecutar los benchmarks en un entorno controlado.

El flujo de trabajo típico es:
\begin{verbatim}
docker build -t ecc-thesis .
docker run -it ecc-thesis /bin/bash
# Dentro del contenedor:
bash run_benchmarks.sh 20
\end{verbatim}

Esto asegura que dos ejecuciones en la misma máquina produzcan resultados comparables, independientemente de las bibliotecas instaladas en el sistema anfitrion.


% ============================================================================
\section{Generador de números aleatorios}
\label{sec:impl-rng}

\subsection{Diseño del RNG}
\label{subsec:impl-rng-diseno}

El generador de números aleatorios se implementa siguiendo el patron de diseño \textit{Strategy}: la interfaz abstracta \texttt{RNG} define los métodos que todo generador debe ofrecer, y la implementación concreta \texttt{NTLRNG} delega en las funciones de NTL (\texttt{SetSeed}, \texttt{RandomBnd}, \texttt{GenPrime}, etc.).

La característica clave del diseño es el soporte para \emph{semillas fijas}: al inicializar \texttt{NTLRNG} con un valor de semilla concreto (por ejemplo, \texttt{seed = 0}), toda la secuencia de números aleatorios generados es idéntica en cualquier ejecución. Esto permite:

\begin{itemize}
    \item Reproducir exactamente los resultados de un benchmark.
    \item Comparar el rendimiento de dos algoritmos con \emph{los mismos} valores aleatorios.
    \item Depurar errores en operaciones criptográficas con datos predecibles.
\end{itemize}

En modo ``aleatorio'' (\texttt{-s random}), la semilla se deriva del timestamp del sistema, proporcionando aleatoriedad real para cada ejecución.


\subsection{Validación estadística}
\label{subsec:impl-rng-validacion}

Para verificar que el RNG produce secuencias con propiedades estadísticas adecuadas para uso criptográfico, se implementó un programa de análisis dedicado (\texttt{rng\_analysis.cpp}) y un conjunto de tests estadísticos en Python (\texttt{analyze\_randomness.py}). Los tests realizados fueron:

\begin{description}
    \item[Test de uniformidad (Chi-cuadrado):] Verifica que la distribución de los valores generados se ajusta a una distribución uniforme.
    \item[Test de Kolmogorov-Smirnov:] Compara la distribución empirica con la teórica uniforme en $[0, 1]$.
    \item[Test de rachas (\textit{runs}):] Verifica que las secuencias de bits consecutivos iguales tienen la longitud esperada bajo aleatoriedad.
    \item[Autocorrelación:] Mide la correlación entre valores consecutivos (debe ser cercana a cero).
    \item[Análisis de entropia:] Calcula la entropia de Shannon de la distribución y la compara con el máximo teórico.
\end{description}

Los resultados de la validación, incluyendo histogramas, gráficas de autocorrelación y resúmenes estadísticos, se generan automáticamente mediante el script \texttt{run\_rng\_analysis.sh}. Todos los tests pasaron satisfactoriamente con $p$-values superiores a 0.05, confirmando que el RNG de NTL produce secuencias estadísticamente indistinguibles de secuencias verdaderamente aleatorias bajo los criterios evaluados.


% ============================================================================
\section{Motor de benchmarking}
\label{sec:impl-benchmark}

\subsection{Metodología de medición}
\label{subsec:impl-metodologia}

Cada benchmark sigue un protocolo estandarizado:

\begin{enumerate}
    \item Se ejecutan 3 iteraciones de calentamiento (\textit{warm-up}) que se descartan, para asegurar que las cachés de CPU estén pobladas y que cualquier inicialización perezosa de NTL se haya completado.
    \item Se ejecutan $N$ iteraciones medidas (configurable con \texttt{-i}).
    \item Para cada iteración, se mide el tiempo transcurrido con \texttt{std::chrono::high\_resolution\_clock} con precisión de microsegundos.
    \item Se calculan las siguientes estadísticas: media, mediana, desviación estándar, mínimo, máximo, y percentiles P5 y P95.
\end{enumerate}

El núcleo del sistema de medición es la función genérica \texttt{run\_benchmark}, implementada como un template de C++ que acepta cualquier operación como argumento lambda:

\begin{verbatim}
template<typename Func>
BenchmarkResult run_benchmark(
    const string& label, Func operation, int iterations);
\end{verbatim}

Esta abstracción permite medir cualquier operación criptográfica (generación de claves, cifrado, firma, etc.) con la misma infraestructura de medición, eliminando la posibilidad de inconsistencias entre benchmarks.


\subsection{Operaciones medidas}
\label{subsec:impl-operaciones}

El cuadro~\ref{tab:operaciones} resume las operaciones medidas para cada algoritmo.

\begin{table}[htbp]
\centering
\small
\begin{tabular}{llp{6.5cm}}
\hline
\textbf{Algoritmo} & \textbf{Operación} & \textbf{Descripción} \\
\hline
RSA              & Keygen           & Generación de claves (1024, 2048, 3072, 4096 bits) \\
                 & Encrypt          & Cifrado $c = m^e \bmod n$ \\
                 & Decrypt          & Descifrado directo $m = c^d \bmod n$ \\
                 & Decrypt CRT      & Descifrado con Teorema Chino del Resto \\
                 & Sign             & Firma $s = h^d \bmod n$ \\
                 & Verify           & Verificación $h = s^e \bmod n$ \\
\hline
ECC (afín)       & Keygen           & Generación $Q = dG$ (P-256, P-384, secp256k1) \\
                 & Scalar mult      & Multiplicación escalar $kP$ \\
                 & ECDH             & Secreto compartido $S = d_A Q_B$ \\
                 & ECDSA sign       & Firma con nonce $k$ \\
                 & ECDSA verify     & Verificación $u_1G + u_2Q$ \\
                 & Point add/double & Operaciones elementales de grupo \\
\hline
ECC (Jacobiano)  & Keygen           & Igual que afín, con \texttt{use\_jacobian=true} \\
                 & ECDH             & Secreto compartido con aritmética Jacobiana \\
                 & ECDSA sign/verify & Firma y verificación con aritmética Jacobiana \\
\hline
ECC (binario)    & Keygen           & Generación sobre $\mathbb{F}_{2^m}$ (5 curvas SEC~2) \\
                 & Scalar mult      & Multiplicación escalar en cuerpo binario \\
                 & ECDH             & Secreto compartido sobre $\mathbb{F}_{2^m}$ \\
\hline
\end{tabular}
\caption{Operaciones criptográficas medidas en el benchmark.}
\label{tab:operaciones}
\end{table}

Todas las comparaciones RSA-ECC se realizan a niveles de seguridad equivalentes según NIST SP~800-57: RSA-3072 frente a P-256/secp256k1 (128 bits de seguridad) y RSA-4096 frente a P-384 (192 bits de seguridad). 


\subsection{Pipeline de visualización}
\label{subsec:impl-visualizacion}

Los resultados del benchmark se procesan mediante dos scripts de Python:

\begin{description}
    \item[\texttt{visualize\_benchmarks.py}] Lee los ficheros CSV generados por el motor de benchmarking y produce 11 gráficas individuales: comparaciones de generación de claves, operaciones de firma y verificación, speedup de ECDH, escalado por tamaño de clave, comparación afín vs.\ Jacobiano, y rendimiento en cuerpos binarios. Utiliza una paleta de cuatro colores consistente (azul/rojo/verde/morado) para distinguir los tipos de algoritmo.
    
    \item[\texttt{create\_collages.py}] Combina las gráficas individuales en 3 collages temáticos correspondientes a los tres ejes de comparación, facilitando su inclusión en la memoria y en presentaciones. Utiliza PIL (Pillow) para la composición de imágenes.
\end{description}

El pipeline completo se orquesta mediante \texttt{run\_benchmarks.sh}, que automatiza la compilación, ejecución de benchmarks, verificación de datos y generación de visualizaciones en un único comando:

\begin{verbatim}
bash run_benchmarks.sh 20   # 20 iteraciones por benchmark
\end{verbatim}

El script genera ficheros con timestamp (evitando sobreescrituras) y crea enlaces simbolicos \texttt{summary\_latest.csv} y \texttt{raw\_latest.csv} para facilitar el acceso a los resultados más recientes.


% ============================================================================
\section{Curvas implementadas}
\label{sec:impl-curvas}

\subsection{\texorpdfstring{Curvas sobre cuerpos primos $\mathbb{F}_p$}{Curvas sobre cuerpos primos Fp}}
\label{subsec:impl-curvas-primas}

Se implementaron tres curvas estándar, con parámetros tomados directamente de FIPS~186-4 y SEC~2:

\begin{table}[htbp]
\centering
\begin{tabular}{lcccc}
\hline
\textbf{Curva} & \textbf{Bits} & \textbf{$a$} & \textbf{$h$} & \textbf{Seguridad} \\
\hline
NIST P-256 (secp256r1) & 256 & $-3$ & 1 & $\sim$128 bits \\
NIST P-384             & 384 & $-3$ & 1 & $\sim$192 bits \\
secp256k1 (Bitcoin)    & 256 & $0$  & 1 & $\sim$128 bits \\
\hline
\end{tabular}
\caption{Curvas sobre cuerpos primos implementadas.}
\label{tab:curvas-primas-impl}
\end{table}

Las curvas NIST utilizan $a = -3$, lo que permite una optimización en la formula de duplicación Jacobiana: el termino $3X_1^2 + aZ_1^4$ se simplifica a $3(X_1 - Z_1^2)(X_1 + Z_1^2)$, ahorrando una multiplicación. La curva secp256k1 utiliza $a = 0$, lo que elimina completamente el termino $aZ_1^4$.


\subsection{\texorpdfstring{Curvas sobre cuerpos binarios $\mathbb{F}_{2^m}$}{Curvas sobre cuerpos binarios F2m}}
\label{subsec:impl-curvas-binarias}

Se implementaron cinco curvas del estándar SEC~2~\cite{sec2-v2}:

\begin{table}[htbp]
\centering
\begin{tabular}{lcccl}
\hline
\textbf{Curva} & \textbf{$m$} & \textbf{$h$} & \textbf{Seguridad} & \textbf{Polinomio irreducible} \\
\hline
sect163k1 & 163 & 2 & $\sim$80 bits  & $x^{163} + x^7 + x^6 + x^3 + 1$ \\
sect233k1 & 233 & 4 & $\sim$112 bits & $x^{233} + x^{74} + 1$ \\
sect283k1 & 283 & 4 & $\sim$128 bits & $x^{283} + x^{12} + x^7 + x^5 + 1$ \\
sect233r1 & 233 & 2 & $\sim$112 bits & $x^{233} + x^{74} + 1$ \\
sect283r1 & 283 & 2 & $\sim$128 bits & $x^{283} + x^{12} + x^7 + x^5 + 1$ \\
\hline
\end{tabular}
\caption{Curvas sobre cuerpos binarios implementadas (SEC~2).}
\label{tab:curvas-binarias-impl}
\end{table}

Las curvas con sufijo ``k'' (Koblitz) tienen $a \in \{0, 1\}$, lo que permite optimizaciones basadas en el endomorfismo de Frobenius. Las curvas con sufijo ``r'' (random) utilizan coeficientes pseudoaleatorios verificables. En ambos casos, la aritmética del cuerpo se realiza íntegramente con \texttt{NTL::GF2E}, donde la suma es la operación XOR y la multiplicación es el producto de polinomios módulo el polinomio irreducible correspondiente.


% ============================================================================
\section{Limitaciones del prototipo}
\label{sec:impl-limitaciones}

Es importante señalar que esta implementación es un prototipo educativo y de benchmarking, \emph{no} un sistema criptográfico de producción. Las principales limitaciones son:

\begin{enumerate}
    \item \textbf{Resistencia a canales laterales.} El algoritmo \textit{double-and-add} tiene un flujo de ejecución que depende de los bits de la clave privada, lo que lo hace vulnerable a ataques de tiempo, consumo eléctrico y emision electromagnética. Una implementación de producción utilizaría la escalera de Montgomery o algoritmos de tiempo constante.
    
    \item \textbf{Nonces ECDSA.} Los nonces $k$ se generan de forma aleatoria. Una implementación robusta utilizaría nonces deterministas según RFC~6979~\cite{rfc6979} para evitar los riesgos descritos en la sección~\ref{subsec:ecdsa-problemas}.
    
    \item \textbf{Derivación de clave (KDF).} El secreto compartido ECDH se deriva truncando la coordenada~$x$ del punto compartido, en lugar de usar un KDF estándar como HKDF (RFC~5869).
    
    \item \textbf{Esquemas de padding.} RSA se implementa sin padding (textbook RSA). Una implementación segura requeriría OAEP para cifrado y PSS para firmas.
    
    \item \textbf{Validación incompleta de parámetros.} Como se detalla en la sección~\ref{subsec:impl-validacion-params}, no se verifica la primalidad de~$n$, $nG = \mathcal{O}$, ni las condiciones de inmersión y anomalía, confiando en los parámetros estandarizados.
\end{enumerate}

Estas limitaciones no afectan la validez de las comparaciones de rendimiento, que es el objetivo principal del proyecto, pero deben tenerse en cuenta si se considerase reutilizar el código en un contexto de seguridad real.